\chapter*{Introdução}
\markboth{Introdução}{Introdução}

Este livro nasceu de uma convicção simples: \textbf{Jesus reina}, e portanto é possível viver com \textbf{prudência, esperança e fidelidade} em meio a sistemas cada vez mais digitais, centralizados e persuasivos. 
O que chamamos aqui de \textit{apocalipse tecnológico} não é um ``novo evangelho'', mas a \textbf{convergência} entre velhos padrões bíblicos de idolatria e \textbf{ferramentas} modernas que amplificam \emph{narrativas}, \emph{identidades} e \emph{controles}.

\section*{A grande ideia: NIC}
A espinha dorsal deste livro é o \textbf{framework NIC} --- \textit{Narrativa, Identidade e Controle}. 
Ele aparece em boxes, exercícios e checklists ao longo dos capítulos:
\begin{nicbox}
\begin{itemize}[leftmargin=*,noitemsep]
  \item \textbf{Narrativa} (\emph{o que é verdade?}) --- quem conta a história e como ela é legitimada.
  \item \textbf{Identidade} (\emph{quem eu sou?}) --- selos, marcas, nomes, pertencimento.
  \item \textbf{Controle} (\emph{o que posso fazer?}) --- portas de compra e venda, permissões, coerções.
\end{itemize}
\end{nicbox}
Quando o leitor aprender a enxergar NIC nos textos bíblicos \emph{e} nos sistemas digitais, ganhará uma bússola simples para tempos confusos.

\section*{Para quem escrevo}
Para cristãos comuns, líderes e profissionais de tecnologia que desejam discernir os sinais do nosso tempo sem pânico nem ingenuidade. \textbf{Não} peço que você concorde com todas as minhas leituras, apenas que verifique \textit{com a Bíblia aberta} e aplique com amor à sua família e igreja.

\section*{Como usar este livro}
Cada capítulo segue uma arquitetura constante: \textit{vinheta de abertura} (história curta), \textit{mapa bíblico enxuto}, \textit{NIC aplicado}, \textit{ferramenta/checklist} e \textit{perguntas para grupo}. Feche cada capítulo com uma \textbf{pequena prática} na semana: são passos modestos que, somados, formam um Plano de 100 Dias (reunido no apêndice).

\section*{Sobre apócrifos e debates escatológicos}
Cito eventualmente literatura judaica do período intertestamentário (p.\,ex., 1~Enoque) como \textit{contexto histórico}, não como autoridade canônica. Quanto às posições escatológicas (premilenismo, amilenismo, pós-milenismo), priorizo o que é \textbf{essencial} e apologio por caridade no que é \textbf{secundário}. O leitor encontrará um mapa honesto no final do livro.

\section*{Convites}
\begin{checklist}[title={Antes de começar}]
\begin{itemize}[leftmargin=*,noitemsep]
  \item Separe uma Bíblia de papel e um caderno. Desligue notificações durante a leitura.
  \item Comprometa-se com \textbf{30 minutos semanais} para ler Daniel 7, Mateus 24, 2~Tessalonicenses~2 e Apocalipse 1.
  \item Crie no caderno três abas: \textbf{N} (Narrativa), \textbf{I} (Identidade), \textbf{C} (Controle).
  \item Combine com alguém de sua igreja de lerem juntos um capítulo por semana.
\end{itemize}
\end{checklist}

Seja bem-vindo. O objetivo não é “vencer discussões na internet”, mas \textbf{cuidar de gente} sob o senhorio de Cristo. Quando a cidade apagar, que sua casa brilhe.
